% GENERATED FILE. Edit canonical lessons, then run npm run generate:books.
% canonical-source-hash: fnv1a64:716cc8eab80c3849
% canonical-lessons: MR-C156-dakhavne, MR-C156-vajavne, MR-C156-oradne, MR-C156-cavne, MR-C156-gilne

\chapter{To show, To play, To shout, To bite, To swallow}
\label{ch:mr-to-show-to-play-to-shout-to-bite-to-swallow}
\hlchaptermodality{156}

\begin{chapteropening}
\textbf{By the end of this chapter:} \emph{I can say to show, to play, to shout, to bite and to swallow.}

Complete the last lesson of chapter 156: i can say to show, to play, to shout, to bite and to swallow.
\end{chapteropening}

\section[dākhavṇe]{\mr{दाखवणे} (dākhavṇe) — to show}

\label{lesson:MR-C156-dakhavne}



\begin{quote}
Before the new one: say the Marathi for to paint, then the Marathi for to take out.
\end{quote}

\subsection*{You'll want to know: \mr{दाखवणे}}
\textbf{\mr{दाखवणे}} — \emph{dākhavṇe} — \textquotedblleft{}to show\textquotedblright{}.

\begin{grammarlens}[title={What you've built}]
One more word for this chapter.
\end{grammarlens}

\subsection*{Your turn}
\begin{itemize}
\raggedright
  \item \emph{Say it:} \emph{dākhavṇe}
  \item \emph{Say it:} \emph{dākhavṇe}, once more
  \item \emph{Say it:} say \emph{dākhavṇe}
  \item \emph{Recall:} say the Marathi for to paint, then the Marathi for to take out, then say \emph{dākhavṇe} again
\end{itemize}

\begin{tcolorbox}[breakable,colback=teal!4,colframe=teal!35!black,title={Before you move on}]
What does \mr{दाखवणे} mean? (\textquotedblleft{}To show\textquotedblright{}.) Say it once more.
\end{tcolorbox}
\section[vājavṇe]{\mr{वाजवणे} (vājavṇe) — to play (an instrument)}

\label{lesson:MR-C156-vajavne}



\begin{quote}
Before the new one: say the Marathi for to take out, then the Marathi for to show.
\end{quote}

\subsection*{You'll want to know: \mr{वाजवणे}}
\textbf{\mr{वाजवणे}} — \emph{vājavṇe} — \textquotedblleft{}to play (an instrument)\textquotedblright{}.

\begin{grammarlens}[title={What you've built}]
One more word for this chapter.
\end{grammarlens}

\subsection*{Your turn}
\begin{itemize}
\raggedright
  \item \emph{Say it:} \emph{vājavṇe}
  \item \emph{Say it:} \emph{vājavṇe}, once more
  \item \emph{Say it:} say \emph{vājavṇe}
  \item \emph{Recall:} say the Marathi for to take out, then the Marathi for to show, then say \emph{vājavṇe} again
\end{itemize}

\begin{tcolorbox}[breakable,colback=teal!4,colframe=teal!35!black,title={Before you move on}]
What does \mr{वाजवणे} mean? (\textquotedblleft{}To play (an instrument)\textquotedblright{}.) Say it once more.
\end{tcolorbox}
\section[oraḍṇe]{\mr{ओरडणे} (oraḍṇe) — to shout}

\label{lesson:MR-C156-oradne}



\begin{quote}
Before the new one: say the Marathi for to show, then the Marathi for to play.
\end{quote}

\subsection*{You'll want to know: \mr{ओरडणे}}
\textbf{\mr{ओरडणे}} — \emph{oraḍṇe} — \textquotedblleft{}to shout\textquotedblright{}.

\begin{grammarlens}[title={What you've built}]
One more word for this chapter.
\end{grammarlens}

\subsection*{Your turn}
\begin{itemize}
\raggedright
  \item \emph{Say it:} \emph{oraḍṇe}
  \item \emph{Say it:} \emph{oraḍṇe}, once more
  \item \emph{Say it:} say \emph{oraḍṇe}
  \item \emph{Recall:} say the Marathi for to show, then the Marathi for to play, then say \emph{oraḍṇe} again
\end{itemize}

\begin{tcolorbox}[breakable,colback=teal!4,colframe=teal!35!black,title={Before you move on}]
What does \mr{ओरडणे} mean? (\textquotedblleft{}To shout\textquotedblright{}.) Say it once more.
\end{tcolorbox}
\section[cāvṇe]{\mr{चावणे} (cāvṇe) — to bite, to chew}

\label{lesson:MR-C156-cavne}



\begin{quote}
Before the new one: say the Marathi for to play, then the Marathi for to shout.
\end{quote}

\subsection*{You'll want to know: \mr{चावणे}}
\textbf{\mr{चावणे}} — \emph{cāvṇe} — \textquotedblleft{}to bite, to chew\textquotedblright{}.

\begin{grammarlens}[title={What you've built}]
One more word for this chapter.
\end{grammarlens}

\subsection*{Your turn}
\begin{itemize}
\raggedright
  \item \emph{Say it:} \emph{cāvṇe}
  \item \emph{Say it:} \emph{cāvṇe}, once more
  \item \emph{Say it:} say \emph{cāvṇe}
  \item \emph{Recall:} say the Marathi for to play, then the Marathi for to shout, then say \emph{cāvṇe} again
\end{itemize}

\begin{tcolorbox}[breakable,colback=teal!4,colframe=teal!35!black,title={Before you move on}]
What does \mr{चावणे} mean? (\textquotedblleft{}To bite, to chew\textquotedblright{}.) Say it once more.
\end{tcolorbox}
\section[giḷṇe]{\mr{गिळणे} (giḷṇe) — to swallow}

\label{lesson:MR-C156-gilne}



\begin{quote}
Before the new one: say the Marathi for to shout, then the Marathi for to bite.
\end{quote}

\subsection*{You'll want to know: \mr{गिळणे}}
\textbf{\mr{गिळणे}} — \emph{giḷṇe} — \textquotedblleft{}to swallow\textquotedblright{}.

\begin{grammarlens}[title={What you've built}]
One more word for this chapter.
\end{grammarlens}

\subsection*{Your turn}
\begin{itemize}
\raggedright
  \item \emph{Say it:} \emph{giḷṇe}
  \item \emph{Say it:} \emph{giḷṇe}, once more
  \item \emph{Say it:} say \emph{giḷṇe}
  \item \emph{Recall:} say the Marathi for to shout, then the Marathi for to bite, then say \emph{giḷṇe} again
\end{itemize}

\begin{tcolorbox}[breakable,colback=teal!4,colframe=teal!35!black,title={Before you move on}]
What does \mr{गिळणे} mean? (\textquotedblleft{}To swallow\textquotedblright{}.) Say it once more.
\end{tcolorbox}
